\SourceLeaf{419}
\[
\dfrac{\operatorname{sin\,lemn}\dfrac{A-B}{1+i}}
{\operatorname{sin\,lemn}\dfrac{A+B}{1+i}}
=\dfrac{\operatorname{sin\,lemn}A-\operatorname{sin\,lemn}B}
{\operatorname{sin\,lemn}A+\operatorname{sin\,lemn}B}
\]
\[
pp=QQ-PP,\qquad qq=QQ+PP
\]
\[
\begin{aligned}P(a-b)\, .\,Q(a+b)&=Pa\, .\,Qa\, .\,\surd(Qb^4-Pb^4)\\[5pt]&\quad{}-Pb\, .\,Qb\, .\,\surd(Qa^4-Pa^4)\end{aligned}
\]
\[
\begin{aligned}P(a-b)\, .\,P(a+b)&=Pa^2\, .\,Qb^2-Qa^2\, .\,Pb^2\end{aligned}
\]
\[
\begin{aligned}Q(a-b)\, .\,Q(a+b)&=Pa^2\, .\,Pb^2+Qa^2\, .\,Qb^2\end{aligned}
\]
\[
\begin{aligned}
q(a-b)\, .\,q(a+b)&=pa^2\, .\,Pb^2+qa^2\, .\,Qb^2\\
&=pb^2\, .\,Pa^2+qb^2\, .\,Qa^2
\end{aligned}
\]
\[
\begin{aligned}q(a-b)\, .\,Q(a+b)&=qa\, .\,qb\, .\,Qa\, .\,Qb\\[5pt]&\quad{}-pa\, .\,pb:Pa\, .\,Pb\end{aligned}
\]
\[
\begin{aligned}q(a-b)\, .\,P(a+b)&=qa\, .\,pb\, .\,Pa\, .\,Qb\\[5pt]&\quad{}+pa\, .\,qb\, .\,Qa\, .\,Pb\end{aligned}
\]
\[
\begin{aligned}
P\varphi&=P\\[4pt]
Q\varphi&=Q\\[4pt]
p\varphi&=\surd(QQ-PP)\\[4pt]
q\varphi&=\surd(QQ+PP)
\end{aligned}
\]
\[
\begin{aligned}
P2\varphi&=2PQ\surd(Q^4-P^4)\\[4pt]
Q2\varphi&=Q^4+P^4\\[4pt]
p2\varphi&=Q^4-2QQPP-P^4\\[4pt]
q2\varphi&=Q^4+2QQPP+P^4
\end{aligned}
\]
\[
\begin{aligned}P3\varphi&=3Q^8P-6Q^4P^5-P^9\end{aligned}
\]
\[
\begin{aligned}Q3\varphi&=Q^9+6Q^5P^4-3QP^8\end{aligned}
\]
\[
\begin{aligned}p3\varphi&=\surd(QQ-PP)\, .\,(Q^8-4Q^6PP\\[5pt]&\quad{}-6Q^4P^4-4QQP^6+P^8)\end{aligned}
\]
\[
\begin{aligned}q3\varphi&=\surd(QQ+PP)\, .\,(Q^4+4Q^6PP\\[5pt]&\quad{}-6Q^4P^4+4QQP^6+P^8)\end{aligned}
\]
\[
\begin{aligned}P4\varphi&=4PQ\surd(Q^4-P^4)\\[5pt]&\quad{}\, .\,(Q^{12}-5Q^8P^4-5Q^4P^8+P^{12})\end{aligned}
\]
\[
\begin{aligned}Q4\varphi&=Q^{16}+20Q^{12}P^4-26Q^8P^8\\[5pt]&\quad{}+20Q^4P^{12}+P^{16}\end{aligned}
\]
\[
\begin{aligned}p4\varphi&=Q^{16}-8Q^{14}PP-12Q^{12}P^4\\[5pt]&\quad{}-8Q^{10}P^6+38Q^8P^8+8Q^6P^{10}\\[5pt]&\quad{}-12Q^4P^{12}+8QQP^4+P^{16}\end{aligned}
\]
\[
\begin{aligned}q4\varphi&=Q^{16}+8Q^{14}PP-12Q^{12}P^4\\[5pt]&\quad{}+8Q^{10}P^6+38Q^8P^8-8Q^6P^{10}\\[5pt]&\quad{}-12Q^4P^{12}-8QQP^4+P^{16}\end{aligned}
\]


\SourceLeaf{420}
\[
\begin{aligned}
\dfrac{Qn\varphi}{Q\varphi^{nn}}
&=1+\dfrac1{12}(n^4-nn)\left(\dfrac{P\varphi}{Q\varphi}\right)^4
\\[5pt]&\quad{}-\dfrac1{10080}(n^8+70n^6\\[5pt]&\quad{}-371n^4+300nn)\left(\dfrac{P\varphi}{Q\varphi}\right)^8\\[7pt]
&\quad+\dfrac1{19958400}(17n^{12}+165n^{10}\\[5pt]&\quad{}+4191n^8-106865n^6+426492n^4\\[7pt]
&\quad-324000nn)\, .\,\left(\dfrac{P\varphi}{Q\varphi}\right)^{12}\\[7pt]
&=1+\dfrac1{12}nn(n-1)\left(\dfrac{P\varphi}{Q\varphi}\right)^4
\\[5pt]&\quad{}-\dfrac1{10080}nn(nn-1)(nn-4)(nn+75)\left(\dfrac{P\varphi}{Q\varphi}\right)^8\\[7pt]
&\quad+\dfrac1{19958400}nn(nn-1)(nn-4)(nn-9)\\[5pt]&\quad{}(17n^4+403nn+9000)\, .\,\left(\dfrac{P\varphi}{Q\varphi}\right)^{12}
\end{aligned}
\]
\[
\begin{aligned}
Pi\varphi&=iP\varphi\\[4pt]
Qi\varphi&=Q\varphi\\[4pt]
pi\varphi&=q\varphi\\[4pt]
qi\varphi&=p\varphi
\end{aligned}
\]
\[
\begin{aligned}
P(1+i)\varphi&=(1+i)PQ\\[4pt]
Q(1+i)\varphi&=\surd(Q^4-P^4)=pq\\[4pt]
p(1+i)\varphi&=QQ-iPP\\[4pt]
q(1+i)\varphi&=QQ+iPP
\end{aligned}
\]
\[
\begin{aligned}P(2+i)\varphi&=(2+i)PQ^4-iP^5\end{aligned}
\]
\[
\begin{aligned}Q(2+i)\varphi&=Q^5-(1-2i)QP^4\end{aligned}
\]
\[
\begin{aligned}p(2+i)\varphi&=\surd(QQ+PP)\\[5pt]&\quad{}\, .\,\{Q^4-(2+2i)QQPP+P^4\}\end{aligned}
\]
\[
\begin{aligned}q(2+i)\varphi&=\surd(QQ-PP)\\[5pt]&\quad{}\, .\,\{Q^4+(2+2i)QQPP+P^4\}\end{aligned}
\]
\[
\begin{aligned}P(3+i)\varphi&=(3+i)PQ^3-(2+6i)P^5Q^5\\[5pt]&\quad{}+(3+i)P^9Q\end{aligned}
\]
\[
\begin{aligned}Q(3+i)\varphi&=\surd(Q^4-P^4)\\[5pt]&\quad{}\, .\,\{Q^8+(2+8i)P^4Q^4+P^8\}\end{aligned}
\]
\[
\begin{aligned}p(3+i)\varphi&=Q^{10}-(4-3i)Q^8PP\\[5pt]&\quad{}-(2-4i)Q^6P^4+(4+2i)Q^4P^6\\[5pt]&\quad{}-(3+4i)QQP^8-iP^{10}\end{aligned}
\]
\[
\begin{aligned}q(3+i)\varphi&=Q^{10}+(4-3i)Q^8PP\\[5pt]&\quad{}-(2-4i)Q^6P^4-(4+2i)Q^4P^6\\[5pt]&\quad{}-(3+4i)QQP^8+iP^{10}\end{aligned}
\]
\[
\begin{aligned}P(1+2i)\varphi&=(1+2i)PQ^4-P^5\end{aligned}
\]
\[
\begin{aligned}Q(1+2i)\varphi&=Q^5-(1+2i)QP^4\end{aligned}
\]
\[
\begin{aligned}P(1+3i)\varphi&=(1+3i)PQ^9-(6+2i)P^5Q^5\\[5pt]&\quad{}+(1+3i)P^9Q\end{aligned}
\]
\[
\begin{aligned}Q(1+3i)\varphi&=Q^{10}\mathord{.}\mathord{.}\end{aligned}
\]
\[
\begin{aligned}P(1+4i)\varphi&=(1+4i)PQ^{16}-(20+12i)P^5Q^{12}\\[5pt]&\quad{}-(10-28i)P^9Q^8\\[5pt]&\quad{}+(12-20i)P^{13}Q^4+P^{17}\end{aligned}
\]
\[
\begin{aligned}Q(1+4i)\varphi&=\end{aligned}
\]

\SourceLeaf{421}
\[
\begin{aligned}
P(1+ni)\varphi
 &= (1+in)PQ^{nn}
 \\[5pt]&\quad{}-\left(\dfrac{nn.nn-1}{12}+\dfrac{n.nn-1.nn-4}{60}i\right)P^5Q^{nn-4}+{}\\[7pt]
 &= (1+in)PQ^{nn}\\[5pt]&\quad{}-\dfrac{n.nn-1.1+in.n-4i}{60}P^5Q^{nn-4}
\end{aligned}
\]
Unter den Zahlen $y=x,\ 2x+1,\ 2x+i,\ 2x-1,\ 2x-i$ ist immer wenigstens eine (oder drei oder alle), die die Auflösung der Congruenz $1-y^4\equiv zz$ nach irgend einem Modulus möglich macht.
\[
\slemn X=x,\qquad \slemn Y=y,\qquad \slemn Z=z
\]
\[
\begin{gathered}\slemn(X+Y+Z)=\\[7pt]
\dfrac{\begin{aligned}
 &-2xyz(xx+yy+zz-xxyyzz)\\[4pt]&\quad{}+x\surd(1-y^4).\surd(1-z^4)\\[3pt]&\quad{}.(1-yyzz+xxzz+xxyy)\\
 &+y\surd(1-x^4).\surd(1-z^4)\\[3pt]&\quad{}.(1+yyzz-xxzz+xxyy)\\
 &+z\surd(1-x^4).\surd(1-y^4)\\[3pt]&\quad{}.(1+yyzz+xxzz-xxyy)
\end{aligned}}{\begin{gathered}1+2yyzz+2xxzz+2xxyy+y^4z^4\\[4pt]{}+x^4z^4+x^4y^4-2x^4yyzz\\[4pt]{}-2xxy^4zz-2xxyyz^4\end{gathered}}\end{gathered}
\]

Ist
\[
\begin{array}{rll}
\varphi+\varphi'+\varphi''=0,\\[5pt]&\slemn\varphi=x&\clemn\varphi=y\\[4pt]
&\slemn\varphi'=x'&\clemn\varphi'=y'\\[4pt]
&\slemn\varphi''=x''&\clemn\varphi'=y''
\end{array}
\]
so ist
\[
\begin{aligned}x(y-y'y'')&=x'(y'-yy'')=x''(y''-yy'),\\[7pt]&\qquad
y''=\dfrac{xy-x'y'}{xy'-x'y}\end{aligned}
\]
\begin{center}
\small
\begin{tabular}{r|r|l}
\multicolumn{1}{c}{$Q$}&\multicolumn{1}{c}{}&\multicolumn{1}{c}{}\\
$-1+2i$&$5$&$1+(-1+2i)x^4$\\[9pt]
$-3$&$9$&$1+6x^4-3x^8$\\[9pt]
$+3+2i$&$13$&$\begin{aligned}[t]1+(-11+10i)x^4\\[3pt]{}+(7-4i)x^8+(3+2i)x^{12}\end{aligned}$\\[9pt]
$+1+4i$&$17$&$\begin{aligned}[t]1+(+12-20i)x^4\\[3pt]{}+(-10+28i)x^8-(20+12i)x^{12}\\[3pt]{}+(1+4i)x^{16}\end{aligned}$\\[9pt]
$5$&$25$&$\begin{aligned}[t]1+50x^4-125x^8\\[3pt]{}+300x^{12}-105x^{16}\\[3pt]{}-62x^{20}+5x^{24}\end{aligned}$
\end{tabular}
\end{center}

für \hfill $x^4=+1$\quad und\quad $x^4=-1$\hfill\mbox{}

werden diese Functionen respective den Quadraten und Würfeln
\[
\begin{array}{rccccc}
\text{von}&1-i&-2&-2-2i&-4i&+8\\
\text{gleich für }M=&5&9&13&17&25
\end{array}
\]
\[
\begin{aligned}
Q(a+bi)\varphi={}&Q^{aa+bb}+\dfrac1{12}\{(a+bi)^4-(aa+bb)\}Q^{aa+bb-4}P^4\\[7pt]
&-\dfrac1{10080}\{(a+bi)^8+70(a+bi)^4(aa+bb)\\[5pt]&\quad{}-35(aa+bb)^2-336(a+bi)^4\\[5pt]&\quad{}+300(aa+bb)\}.\\[7pt]
&\quad.Q^{aa+bb-8}P^8\mathord{.}\mathord{.}\mathord{.}
\end{aligned}
\]
\par\begin{center}\rule{.22\linewidth}{.3pt}\end{center}
